\documentclass[11pt]{article}
\usepackage[a4paper,margin=23mm]{geometry}
\usepackage{amsmath,amssymb,amsthm,mathtools,bm,mathrsfs}
\usepackage{slashed}
\usepackage{booktabs,longtable,array}
\usepackage{enumitem}
\usepackage{xcolor}
\usepackage{microtype}
\usepackage[hidelinks]{hyperref}

\setlength{\emergencystretch}{3em}
\newtheorem{theorem}{Theorem}[section]
\newtheorem{proposition}[theorem]{Proposition}
\newtheorem{lemma}[theorem]{Lemma}
\newtheorem{corollary}[theorem]{Corollary}
\theoremstyle{definition}
\newtheorem{definition}[theorem]{Definition}
\newtheorem{gate}[theorem]{Fail-closed gate}
\theoremstyle{remark}
\newtheorem{remark}[theorem]{Remark}

\newcommand{\ii}{\mathrm i}
\newcommand{\dd}{\mathrm d}
\newcommand{\CC}{\mathbb C}
\newcommand{\RR}{\mathbb R}
\newcommand{\ZZ}{\mathbb Z}
\newcommand{\CP}{\mathbb {CP}}
\newcommand{\cO}{\mathcal O}
\newcommand{\cL}{\mathcal L}
\newcommand{\cM}{\mathcal M}
\newcommand{\cC}{\mathcal C}
\newcommand{\cG}{\mathcal G}
\newcommand{\cE}{\mathcal E}
\newcommand{\cH}{\mathcal H}
\newcommand{\Tr}{\operatorname{Tr}}
\newcommand{\Ind}{\operatorname{Ind}}
\newcommand{\Ker}{\operatorname{Ker}}
\newcommand{\Pf}{\operatorname{Pf}}
\newcommand{\SF}{\operatorname{SF}}
\newcommand{\status}[1]{\textcolor{blue!60!black}{\textsf{#1}}}
\newcolumntype{L}[1]{>{\raggedright\arraybackslash}p{#1}}

\title{AP-E6 Same-Soliton Yukawa--Callias Audit:\\
The Physical Moduli No-Go, a Constituent Dirac Completion,\\
CPT Algebra, and a Typed Local Mixed-WZW Ansatz}
\author{Route F research note}
\date{16 July 2026}

\begin{document}
\maketitle

\begin{abstract}
We test whether the canonical Workline-B relaxed charged two-colour \(B=1\)
field can
support the same physical \(\CP^1\) fermion family required by AP-E4.  The
answer is a sharp no-go for the present model and a correctly typed local
differential-character ansatz; neither is a physics-promotion claim.  The
background is bound to the Workline-B public solver and machine card by the
SHA-256 checksum of the canonical little-endian float64 \((r,F,s)\) array.

First, the internal stabilizer of

\[
 U_F(\bm x)=\cos F(r)+\ii\,\bm\tau\!\cdot\!\widehat{\bm x}\sin F(r)
\]

is only the centre \(\ZZ_2\).  Its pure isospin orbit is therefore
\(SU(2)/\ZZ_2\simeq SO(3)\), not \(\CP^1\).  The same \(SO(3)\) is obtained
from the combined spatial--isospin symmetry after quotienting the diagonal
hedgehog stabilizer; these are not two independent collective-coordinate
sets and neither is a colour-gauge redundancy.  The separate scalar-vacuum
\(\CP^1\) is spatially uniform and has norm proportional to the three-volume,
so it is not a localized normalizable soliton modulus.

Second, we write the minimal declared constituent chiral Yukawa completion
and its static Hamiltonian,
\[
 \mathcal D_E=\gamma^\mu D_\mu
 +M(P_RU_F+P_LU_F^\dagger),\qquad
 H_F=-\ii\alpha^iD_i+\beta M
 \bigl(\cos F+\ii\gamma_5\bm\tau\!\cdot\!\widehat{\bm x}\sin F\bigr).
\]
Writing the Hermitian potential as
\(\Phi_F=\beta M(\cos F+\ii\gamma_5\bm\tau\cdot\widehat{\bm x}\sin F)\),
its asymptotic mass is \(\Phi_\infty=\beta M\), not \(M\bm1\).  Under the
stated domain, bounded-commutator, coercivity, and decay hypotheses, the
ordinary \(S^2_\infty\) Callias positive-eigenbundle is constant and the
static Callias index is zero.  The
four-dimensional interpolation formula
\(\Ind(\partial_\tau+H_\tau)=\SF(H_\tau)=B\) concerns spectral flow and does
not force a zero mode at the final static soliton.  A finite Wilson control
finds no accidental endpoint zero at the benchmark, but does not prove a
continuum spectral gap.

Third, a new two-band obstruction rules out the previously declared mixed
eigenline in the minimal mass representation.  On
\(X=S^2_\infty\times\CP^1\), let \(x,y\) generate the two degree-two
cohomologies.  A globally nondegenerate \(2\times2\) Hermitian mass with one
positive band defines a map \(f:X\to\CP^1\); its eigenline obeys
\(c_1^2=f^*(u^2)=0\).  But the desired line has
\[
 c_1=x+2y,\qquad c_1^2=4xy\ne0.
\]
Thus a nontrivial ambient bundle, at least three bands, or another microscopic
mechanism is necessary.  The numerical \(\cO(2)\) Veronese line is retained
only as a spectator control.

Finally, define first the universal differential character
\[
 \widehat\Xi_5
 =2\,\operatorname{pr}_U^*\widehat\omega_3
 \smile\operatorname{pr}_s^*\widehat c_1(\mathsf H)
 \in\widehat H^5(SU(2)\times\CP^1;\ZZ),
\]
then pull it back by the evaluation map on
\(\Sigma_3\times\cC_{\rm restr}\) and fiber-integrate to a degree-two
character.  This is a well-typed local ansatz.  Constant-loop and centre
phase controls do not prove descent on the restricted gauge stack: all
components of \(\operatorname{Map}(\Sigma_3,U(1))\), an equivariant cocycle,
and vertical holonomy remain untreated.  The odd/even centre test is only a
necessary control, not an exactly-two theorem.  Consequently every descent,
same-soliton composition, and portal gate remains false.
\end{abstract}

\tableofcontents

\section{Question, evidence levels, and theorem-level answer}

The previous AP-E4 note supplied a conditional families-Callias template.
The present task is stronger: start from the \emph{same} nonlinear \(B=1\)
configuration published by charged two-colour Workline-B, determine its true
symmetry orbit, and only then ask for a fermion family, determinant line, CPT
map, and WZW line.  We use four evidence labels.

\begin{description}
\item[Same-background result.] A statement computed from the canonical
  Workline-B \((F(r),s(r))\) and the fields already present in the charged EFT.
\item[Declared completion.] An explicit new fermion term that is compatible
  with the displayed low-energy symmetries but has not been matched from the
  microscopic charged gauge theory.
\item[Restricted theorem.] A complete construction on a stated open
  subspace, with every required equivariance and globality hypothesis stated.
\item[Spectator control.] A finite matrix or bundle designed to validate an
  index, Berry, or regulator formula but not derived from the mother model.
\end{description}

\begin{theorem}[Strongest defensible AP-E6 result]
For the present \(B=1\) hedgehog, the localized rotational/isospin orbit is
\(SO(3)\), while the scalar \(\CP^1\) vacuum orientation is nonnormalizable in
infinite volume.  Hence there is no same-soliton physical \(\CP^1\) on which
an \(\cO(2)\) families determinant line or its \(\cO(-4)\) CPT partner can be
defined.  The minimal constituent chiral Yukawa operator is Fredholm at the
essential-spectrum level but has trivial ordinary Callias boundary class.
Its four-dimensional spectral flow is logically distinct from a static
kernel.  A universal degree-five differential character and its degree-two
family transgression are correctly typed, but restricted gauge-stack descent
is not proved.  In particular, tested constant loops and the faithful centre
do not exhaust the gauge group.  Therefore
\[
 \begin{gathered}
 \texttt{same\_physical\_moduli\_space\_derived=false},\qquad
 \texttt{determinant\_line\_o2\_derived=false},\\
 \texttt{gauge\_basic\_wzw\_descent\_proven=false},\qquad
 \texttt{lane\_closed=false}.
 \end{gathered}
\]
The narrower distinction is
\[
\begin{gathered}
\texttt{restricted\_higgsed\_sigma\_model\_descent\_proven=false},\\
\texttt{restricted\_local\_differential\_character\_ansatz\_typed=true}.
\end{gathered}
\]
\end{theorem}

Recent Skyrmion--fermion computations explicitly find a level that flows
through zero as a Yukawa parameter is varied; they do not turn every static
endpoint into a protected zero mode~\cite{DzhunushalievEtAl2024APE6}.  Recent
generalized Callias theory likewise emphasizes that the index is a pairing of
the actual potential class with the Dirac class and can be localized to an
appropriate hypersurface~\cite{VanDenDungen2025APE6}.  These results motivate
the separation enforced below.

\section{The actual B=1 moduli: stabilizer before quantization}

\subsection{Relaxed hedgehog and its baryon number}

Workline-B solves the coupled nonlinear \((F,s)\) boundary-value problem and
publishes a canonical sampling with box radius (R=16) and (4097) points.
Its pion component defines
\begin{equation}
 U_F(\bm x)=\cos F(r)+\ii\tau_i\widehat x_i\sin F(r),
 \qquad F(0)=\pi,\quad F(\infty)=0.
 \label{eq:hedgehog}
\end{equation}
With the convention
\begin{equation}
 B=-\frac{1}{24\pi^2}\int_{\RR^3}
 \Tr\bigl(U_F^{-1}\dd U_F\bigr)^3,
 \label{eq:baryon}
\end{equation}
the canonical numerical solution gives
\(B=0.999999999686391\) and
\(s(0)=0.917128964607058\).  Neither topology nor spherical symmetry
identifies the moduli space with \(S^2\).

The verifier calls exactly
\begin{center}
\small
\nolinkurl{scan_ap_e6_relaxed_b1_multigrid_hessian}\\[-1pt]
\nolinkurl{.solve_relaxed_profile(box=16,sample_points=4097)}
\end{center}
and forms the SHA-256 digest of the little-endian float64 columns
\((r,F,s)\).  It hard-asserts equality with the digest stored in the
Workline-B machine card:
\begin{center}
\small\nolinkurl{81104f59a5f3f4337739fc2a217cafc8da91581c9f1fb40b4fdba6ce0f948d59}.
\end{center}
The full Workline-B script and card hashes are recorded in
Table~\ref{tab:provenance} and in the JSON source manifest.  Thus
``same soliton'' means byte-identical canonical \((r,F,s)\) data, not an
independent AP-E3 solve.

\subsection{Pure internal stabilizer}

Let \(A\in SU(2)_V\) act by
\(U_F(\bm x)\mapsto A U_F(\bm x)A^{-1}\).
Choose a radius \(r_0\) for which \(\sin F(r_0)\ne0\), which exists for every
nontrivial degree-one profile.  If \(A\) fixes the complete configuration,
then
\begin{equation}
 A(\tau_i\widehat x_i)A^{-1}=\tau_i\widehat x_i
 \quad\text{for all }\widehat{\bm x}\in S^2.
 \label{eq:all-directions}
\end{equation}
Taking successively three independent directions implies
\(A\tau_iA^{-1}=\tau_i\) for \(i=1,2,3\).  By irreducibility of the defining
\(SU(2)\) representation, or directly by the Pauli commutators,
\begin{equation}
 \operatorname{Stab}_{SU(2)_V}(U_F)=\{\bm1,-\bm1\}=\ZZ_2.
 \label{eq:pure-stabilizer}
\end{equation}
Thus the pure internal orbit is
\begin{equation}
 \cM_{\rm int}=SU(2)_V/\ZZ_2\simeq SO(3),
 \qquad \dim_\RR\cM_{\rm int}=3.
 \label{eq:so3-orbit}
\end{equation}
The verifier forms the positive Gram matrix
\begin{equation}
 G_{ab}=\sum_{i=1}^3
 \Tr\bigl([\tau_a,\tau_i]^\dagger[\tau_b,\tau_i]\bigr),
 \label{eq:stabilizer-gram}
\end{equation}
whose kernel is the continuous stabilizer Lie algebra.  All three eigenvalues
are strictly positive.

\subsection{Combined spatial--isospin symmetry is the same orbit}

It is essential not to count another three modes.  Let
\((R,A)\in SO(3)_{\rm space}\times SU(2)_V\) act by
\begin{equation}
 U_F(\bm x)\longmapsto A U_F(R^{-1}\bm x)A^{-1}.
 \label{eq:combined-action}
\end{equation}
Writing \(\pi:SU(2)\to SO(3)\) for the double cover, Eq.~\eqref{eq:hedgehog}
is fixed precisely when \(R=\pi(A)\).  Hence the connected stabilizer is the
diagonal lift \(SU(2)_{\rm diag}\), with the appropriate central
identifications, and
\begin{equation}
 \frac{SO(3)_{\rm space}\times SO(3)_V}{SO(3)_{\rm diag}}
 \simeq SO(3).
 \label{eq:combined-orbit}
\end{equation}
Equations~\eqref{eq:so3-orbit} and \eqref{eq:combined-orbit} describe the same
relative spin--isospin orientation of a hedgehog.  They are not two
independent \(SO(3)\)'s.  Moreover, \(U_F\) is a colour singlet.  Microscopic
\(SU(2)_c\) gauge transformations do not quotient this global pion orbit.

\begin{corollary}[No physical \(\CP^1\) from hedgehog rotations]
The relaxed hedgehog has no \(U(1)\) pure-internal stabilizer.  Therefore the
identification \(SU(2)/U(1)=\CP^1\) cannot be made for this solution.
\end{corollary}

\subsection{Why the scalar vacuum sphere is not a soliton modulus}

The charged UV candidate also has a scalar doublet
\(\phi=vz\), \(z^\dagger z=1\), \(z\sim e^{\ii\lambda}z\), so the vacuum
quotient is \(\CP^1\).  If \(z=z(m)\) is spatially uniform, the kinetic norm
of a variation is
\begin{equation}
 \|\delta_m\phi\|^2
 =v^2\int_{\RR^3}\dd^3x\,
 \delta_mz^\dagger(1-zz^\dagger)\delta_mz.
 \label{eq:bulk-norm}
\end{equation}
For any nonzero tangent variation the density approaches a positive
constant, so the norm diverges as the volume.  In a box
\([-L,L]^3\), \(\|\delta_m\phi\|^2\propto8L^3\).  This is a bulk Goldstone
coordinate, not a localized collective coordinate.  Making it slowly
time-dependent produces an inertia that diverges in the infinite-volume
limit.

\begin{gate}[Same physical moduli]
One would have to find a different relaxed charged soliton with a localized
scalar texture and stabilizer \(U(1)\), then prove its horizontal mode is
\(L^2\)-normalizable.  The present solution proves instead
\(\cM_1\simeq SO(3)\) for rotations and a nonnormalizable bulk \(\CP^1\).
\end{gate}

\section{A concrete constituent Yukawa completion on the same profile}

\subsection{Four-dimensional Euclidean operator}

The lowest-dimension chiral constituent coupling compatible with vectorlike
\(U(1)_g\) charge is declared to be
\begin{equation}
 S_{F,Y}=\int\dd^4x\,
 \overline\Psi\left[
 \gamma^\mu(\partial_\mu+\mathcal A_\mu)
 +M\bigl(P_RU_F+P_LU_F^\dagger\bigr)
 \right]\Psi,
 \qquad P_{R,L}=\frac{1\pm\gamma_5}{2}.
 \label{eq:4d-yukawa}
\end{equation}
Here \(\Psi\) carries spin and flavour/isospin; \(\mathcal A_\mu\) denotes the
declared vectorlike gauge connection.  Equation~\eqref{eq:4d-yukawa} is a
standard chiral Yukawa form, but in this project it is a \emph{declared
constituent completion}: the current Workline-B charged EFT supplied no fermion
representation, Yukawa coefficient, or regulator matching.  It must not be
silently promoted to a renormalizable microscopic interaction.

The canonical Workline-B background also contains the relaxed breathing
field \(s(r)\).  Equation~\eqref{eq:4d-yukawa} deliberately declares no
\(s\overline\Psi\Psi\), \(sU_F\), or other breathing-scalar Yukawa coupling.
Therefore \(s\) is \emph{decoupled by this declared completion}; it is not
silently omitted from the imported solution.  Any microscopic matching that
generates an \(s\)-dependent fermion mass changes the operator and requires a
new Callias audit.

Under the Euclidean CPT image
\begin{equation}
 U_F(\bm x)\mapsto U_F^\dagger(-\bm x),\qquad
 \mathcal A_\mu(x)\mapsto\mathcal A_\mu^{\rm CPT}(x),
 \label{eq:cpt-fields}
\end{equation}
assume an antiunitary map \(\Theta\) on the one-particle Hilbert space such
that \(\Theta\ii\Theta^{-1}=-\ii\) and
\begin{equation}
 \mathcal D_E[\Phi^{\rm CPT}]
 =\Theta\,\mathcal D_E[\Phi]\Theta^{-1}.
 \label{eq:cpt-operator}
\end{equation}
This relation is conditional: \(\Theta\) must map the operator domain,
boundary conditions, and spectral cut to their CPT images; the fermion and
regulator representations must be \(\Theta\)-closed; regulator masses must be
real.  It implies determinant conjugation only after those assumptions and a
consistent determinant-line orientation are supplied.  It does not choose a
physical regulator.

\subsection{Static three-dimensional Hamiltonian}

Choose Pauli matrices \(\rho_i\) for the upper/lower Dirac block,
\(\sigma_i\) for spin, and \(\tau_i\) for isospin.  One convenient Hermitian
representation is
\begin{equation}
 \alpha_i=\rho_1\otimes\sigma_i,\qquad
 \beta=\rho_3\otimes\bm1,\qquad
 \gamma_5=\rho_1\otimes\bm1.
 \label{eq:gamma-rep}
\end{equation}
On
\(\cH=L^2(\RR^3,S\otimes\CC^2_\tau)\), define
\begin{equation}
 D_3=-\ii\alpha^iD_i,
 \qquad
 \Phi_F=\beta M\bigl(\cos F
 +\ii\gamma_5\bm\tau\!\cdot\!\widehat{\bm x}\sin F\bigr),
 \qquad H_F=D_3+\Phi_F,
 \label{eq:dirac-potential-split}
\end{equation}
initially on \(C_c^\infty\) and with self-adjoint domain
\(\operatorname{Dom}H_F=H^1\) under the smooth-decaying gauge assumptions
below.  Explicitly,
\begin{align}
 H_F
 &=-\ii\rho_1\otimes\sigma_iD_i\otimes\bm1_\tau
 +M\rho_3\otimes\bm1\otimes\bm1\cos F(r)\notag\\
 &\hspace{18mm}
 -M\rho_2\otimes\bm1\otimes
 \bigl(\bm\tau\!\cdot\!\widehat{\bm x}\bigr)\sin F(r).
 \label{eq:static-hamiltonian}
\end{align}
The second line follows from
\(\ii\beta\gamma_5=-\rho_2\).  Thus \(\Phi_F\) and \(H_F\) are Hermitian.
Notice also that \(\Phi_F^2=M^2\bm1\).

Since \(F(r)\to0\) exponentially in the Workline-B massive-pion profile,
\begin{equation}
 H_F-H_\infty\ \text{is relatively compact},\qquad
 H_\infty=-\ii\alpha^i\partial_i+\beta M.
 \label{eq:relative-compact}
\end{equation}
Consequently
\begin{equation}
 \sigma_{\rm ess}(H_F)=(-\infty,-|M|]\cup[|M|,\infty).
 \label{eq:essential-spectrum}
\end{equation}
This proves a uniform \emph{essential-spectrum} gap \( |M| \) under global
unitary isorotation.  It does not exclude discrete eigenvalues in the gap,
prove \(\inf|\sigma(H_F)|>0\), or establish a uniform Fredholm family over an
unconstructed physical moduli space.

\subsection{Ordinary Callias boundary index is zero}

The self-adjoint Hamiltonian \(H_F\) is not itself the off-diagonal Fredholm
operator whose index is being quoted.  Define
\begin{equation}
 \mathscr C_F=D_3+\ii\Phi_F:
 H^1(\RR^3,S\otimes\CC^2_\tau)\longrightarrow\cH,
 \qquad
 \mathbb B_F=
 \begin{pmatrix}0&\mathscr C_F^\dagger\\
                 \mathscr C_F&0\end{pmatrix},
 \label{eq:callias}
\end{equation}
where \(\mathbb B_F\) acts on the doubled Hilbert space with grading
\(\Gamma=\operatorname{diag}(1,-1)\).  We impose the following sufficient
Callias hypotheses:
\begin{enumerate}[label=(C\arabic*)]
\item \(\RR^3\) is complete and \(D_3\) is self-adjoint on \(H^1\);
\item \(\Phi_F\) is bounded and Hermitian, and
  \([D_3,\Phi_F]\) extends to a bounded operator;
\item outside a compact set \(K\), for some \(\epsilon>0\),
  \begin{equation}
   \Phi_F^2-\lVert[D_3,\Phi_F]\rVert\bm1\ge\epsilon\bm1;
   \label{eq:callias-coercive}
  \end{equation}
\item the gauge field is smooth and decays sufficiently that
  \(H_F-H_\infty\) is relatively compact.
\end{enumerate}
Then \(\mathscr C_F\) is Fredholm, \(\mathbb B_F\) is its graded
self-adjoint double, and the index is determined by the positive
eigenbundle
\(E_+(\Phi_\infty)\to S^2_\infty\)~\cite{Callias1978APE6,VanDenDungen2025APE6}.
For Eq.~\eqref{eq:dirac-potential-split}, the correct boundary potential is
\begin{equation}
 \Phi_\infty=\beta M,
 \qquad E_+(\Phi_\infty)=S^2_\infty\times\CC^4,
 \qquad c_1(E_+)=0,
 \qquad \Ind\mathscr C_F=0.
 \label{eq:actual-callias-zero}
\end{equation}
The rank-four value refers to the displayed Dirac--isospin representation;
only its constancy is topologically relevant.  Equation
\eqref{eq:actual-callias-zero} is asserted only under (C1)--(C4), not for an
arbitrary microscopic gauge background.
The fact that \(U_F:\RR^3\cup\{\infty\}\simeq S^3\to SU(2)\simeq S^3\)
has degree one is a bulk \(\pi_3\) statement.  The ordinary static Callias
boundary theorem sees \(S^2_\infty\) and a constant mass here; it cannot be
replaced by \(B\) without changing the operator problem.

\subsection{Where B does enter: four-dimensional spectral flow}

Let \(U_\tau\) interpolate from the vacuum to \(U_F\), and define
\begin{equation}
 \mathscr D_{\rm flow}=\partial_\tau+H[U_\tau].
 \label{eq:flow-operator}
\end{equation}
With APS boundary conditions, a maintained endpoint gap, and a regulator
whose anomaly normalization is fixed, the standard index/spectral-flow
identity is
\begin{equation}
 \Ind\mathscr D_{\rm flow}=\SF\{H[U_\tau]\}.
 \label{eq:index-sf}
\end{equation}
In the adiabatic chiral problem the anomaly calculation identifies this
spectral flow with the winding \(B\), up to the stated representation and sign
conventions~\cite{APS1975APE6,GoldstoneWilczek1981APE6}.  But spectral flow
counts crossings \emph{along the path}.  It does not assert
\(\Ker H[U_F]\ne0\).  Numerically observed Skyrmion levels indeed cross zero
only at selected coupling values~\cite{DzhunushalievEtAl2024APE6}.

\begin{gate}[Static Callias versus interpolation index]
For the same canonical background, the defensible statements are
\[
 \begin{gathered}
 \text{uniform essential gap true},\\
 \text{ordinary boundary Callias index }0,\\
 \text{conditional interpolation spectral flow }B.
 \end{gathered}
\]
The full endpoint spectral gap and a physical static zero-mode bundle remain
unproved.
\end{gate}

\section{Families determinant line and a two-band obstruction}

\subsection{Conditional push-forward formula}

To compare with the previous template, suppose temporarily that a physical
\(\CP^1\) family \(Y\) existed and a rank-one positive asymptotic eigenbundle
over
\begin{equation}
 X=S^2_\infty\times Y,
 \qquad H^*(X;\ZZ)=\ZZ[x,y]/(x^2,y^2),
 \qquad \int_{S^2_\infty}x=1
 \label{eq:product-base}
\end{equation}
had
\begin{equation}
 c_1(E_+)=p x+q y.
 \label{eq:pq-line}
\end{equation}
Since
\begin{equation}
 \operatorname{ch}(E_+)
 =1+px+qy+pq\,xy,
 \label{eq:chern-character}
\end{equation}
the oriented boundary push-forward gives
\begin{equation}
 \operatorname{rank}\Ind\mathscr C=\epsilon p,
 \qquad
 c_1(\det\Ind\mathscr C)=\epsilon pq\,y.
 \label{eq:family-pushforward}
\end{equation}
Thus \(\epsilon p=1,q=2\) would produce a rank-one \(\cO(2)\) index line.
The Bismut--Freed determinant connection would then refine
Eq.~\eqref{eq:family-pushforward} geometrically~\cite{BismutFreed1986APE6,
BismutFreed1986IIAPE6}.

\subsection{No globally gapped two-band realization of x+2y}

\begin{theorem}[Two-band mixed-line obstruction]
Let \(\Phi:X\to\operatorname{Herm}(2)\) be continuous on
\(X=S^2\times\CP^1\), acting on the trivial bundle \(X\times\CC^2\).  Assume
that its two eigenvalues are globally nondegenerate and separated by zero,
with exactly one positive band.  Then its positive eigenline \(E_+\) obeys
\begin{equation}
 c_1(E_+)^2=0.
 \label{eq:square-zero}
\end{equation}
In particular \(c_1(E_+)=x+2y\) is impossible.
\end{theorem}

\begin{proof}
Subtracting half the trace and dividing by the nonzero eigenvalue separation
deformation-retracts \(\Phi\) to a traceless unit Hermitian matrix
\begin{equation}
 \widehat\Phi=f_a\tau_a,
 \qquad \bm f:X\longrightarrow S^2\simeq\CP^1.
 \label{eq:mass-map}
\end{equation}
The positive eigenline is \(E_+=f^*\mathsf H_+\), where
\(\mathsf H_+\to\CP^1\) is the universal positive eigenline.  If
\(u=c_1(\mathsf H_+)\), then \(u^2=0\) because
\(H^4(\CP^1;\ZZ)=0\).  Naturality gives

\[
 c_1(E_+)^2=f^*(u^2)=0.
\]

On the other hand, using \(x^2=y^2=0\),
\begin{equation}
 (x+2y)^2=4xy,
 \label{eq:mixed-square}
\end{equation}
and \(xy\) generates \(H^4(S^2\times\CP^1;\ZZ)\).  Hence
\(4xy\ne0\), a contradiction.
\end{proof}

The hypotheses matter.  The theorem does not cover a nontrivial ambient
rank-two bundle, a band crossing, more than two bands, or an operator-valued
infinite-dimensional mass.  It does prove that the most economical globally
gapped \(2\times2\) spectator mass cannot realize the desired mixed line.
At least a rank-three ambient construction or nontrivial microscopic bundle
is necessary, but not by itself sufficient.

\subsection{Why the Veronese control is not the missing construction}

For \(z=(z_0,z_1)\in\CC^2\), define
\begin{equation}
 u_2(z)=\bigl(z_0^2,\sqrt2z_0z_1,z_1^2\bigr)^T,
 \qquad
 K(z)=\bm1_3-u_2(z)u_2(z)^\dagger.
 \label{eq:veronese}
\end{equation}
Then \(K u_2=0\), the two nonzero eigenvalues equal one, and
\(u_2(e^{\ii\lambda}z)=e^{2\ii\lambda}u_2(z)\).  The kernel line over
\(\CP^1\) is \(\cO(2)\).  The verifier computes its Berry number \(+2\) using
gauge-invariant Bargmann triangle phases.

This is a rank-three \emph{spectator} control over the \(y\)-sphere only.  It
does not construct the required \(p=1\) dependence over \(S^2_\infty\), does
not couple to Eq.~\eqref{eq:static-hamiltonian}, and is not a Yukawa term of
the charged mother model.

\subsection{The actual orientation space cannot carry free c1=2}

Even before choosing a mass, the physical orbit found above has
\begin{equation}
 H^2(SO(3);\ZZ)\simeq\ZZ_2,
 \qquad H^2(SU(2);\ZZ)=H^2(S^3;\ZZ)=0.
 \label{eq:so3-cohomology}
\end{equation}
Thus an orientation determinant line can at most carry a torsion
\(\ZZ_2\) class on \(SO(3)\); on the spin cover it is topologically trivial.
There is no free Chern number \(+2\) to identify with \(\cO(2)\).  A possible
torsion sign is a Finkelstein--Rubinstein/fermion-measure question and is not
computed here.

\begin{gate}[Families determinant]
The abstract formula~\eqref{eq:family-pushforward} and the Veronese regression
are correct.  The same soliton supplies neither their \(\CP^1\) base nor the
mixed asymptotic eigenbundle.  Therefore
\(\texttt{determinant\_line\_o2\_derived=false}\).
\end{gate}

\section{CPT algebra, Pfaffians, and the conditional anti-canonical map}

\subsection{Pauli--Villars determinant transformation}

Let \(\mathcal D_M[\Phi]\) be Eq.~\eqref{eq:4d-yukawa} and choose a
CPT-covariant Pauli--Villars operator \(\mathcal D_\Lambda[\Phi]\) with the
same kinetic representation.  Define
\begin{equation}
 Z_{\rm PV}[\Phi]
 =\frac{\det\mathcal D_M[\Phi]}{\det\mathcal D_\Lambda[\Phi]}.
 \label{eq:pv-ratio}
\end{equation}
The following are assumptions, not outputs of the charged EFT:
\begin{enumerate}[label=(T\arabic*)]
\item \(\Theta\) is antiunitary, obeys
  \(\Theta\ii\Theta^{-1}=-\ii\), and satisfies
  Eq.~\eqref{eq:cpt-operator};
\item \(\Theta\) maps the operator domain, boundary conditions, and spectral
  cut to their CPT images;
\item the physical and Pauli--Villars representations are closed under
  \(\Theta\), and every regulator mass is real;
\item determinant and Pfaffian orientations are transported consistently.
\end{enumerate}
Under (T1)--(T4), for both masses,
\begin{equation}
 Z_{\rm PV}[\Phi^{\rm CPT}]=Z_{\rm PV}[\Phi]^*.
 \label{eq:det-cpt}
\end{equation}
Equivalently, its eta/phase contribution reverses sign.  A finite spectral
cutoff regression uses real eigenvalues \(\lambda_j\) and
\begin{equation}
 Z_+(\mu,\Lambda)=\prod_j
 \frac{\mu+\ii\lambda_j}{\Lambda+\ii\lambda_j},\qquad
 Z_-=\prod_j
 \frac{\mu-\ii\lambda_j}{\Lambda-\ii\lambda_j}=Z_+^*.
 \label{eq:spectral-regulator}
\end{equation}
This is an algebraic finite spectral-product identity, not a microscopic
regulator construction.  Global
phases of fermion determinants require eta/Dai--Freed data beyond a local
matrix identity~\cite{WittenYonekura2019APE6}.

\subsection{Majorana doubling and Pfaffian orientation}

For an antisymmetric doubled matrix
\begin{equation}
 \mathcal A(D)=
 \begin{pmatrix}0&D\\-D^T&0\end{pmatrix},
 \qquad
 \Pf\mathcal A(D)=(-1)^{n(n-1)/2}\det D,
 \label{eq:pf-doubling}
\end{equation}
Eq.~\eqref{eq:det-cpt} implies the corresponding Pfaffian conjugation after a
fixed orientation.  Whether the orientation can be fixed globally on all
gauge sectors is precisely an anomaly/regulator question and remains open.

\subsection{Anti-canonical O(-4) remains conditional}

If a physical polarized \(\CP^1\) family with \(E_+=\cO(2)\) existed, fixed
Dolbeault polarization and Serre duality would map it to
\begin{equation}
 E_-=K_{\CP^1}\otimes E_+^\vee
 =\cO(-2)\otimes\cO(-2)=\cO(-4).
 \label{eq:anti-canonical}
\end{equation}
Then \((h^0,h^1)=(3,0)\) maps to \((0,3)\).  But
Eq.~\eqref{eq:anti-canonical} is a statement on \(\CP^1\).  The actual
orientation space is \(SO(3)\), and the mother model has not supplied the
polarization or the anti-canonical regulator factor.  Determinant
conjugation alone does not manufacture them.

\begin{gate}[Physical CPT regulator]
The determinant/Pfaffian conjugation law follows conditionally from
(T1)--(T4), and its finite spectral-product regression is algebraically
correct.  No \(\Theta\)-covariant regulator has been derived from the mother
model, and the \(\cO(-4)\) physical line is absent.  Hence
\(\texttt{physical\_cpt\_regulator\_derived=false}\).
\end{gate}

\section{Typed local mixed-WZW character and the unclosed descent problem}

\subsection{Universal character before evaluation}

Restrict first to \(\phi^\dagger\phi=v^2>0\) everywhere.  Write
\(z=\phi/v\in S^3\), with \(z\mapsto e^{\ii\lambda}z\), and let
\begin{equation}
 s=[z]:M_4\longrightarrow\CP^1.
 \label{eq:projective-map}
\end{equation}
Let \(\mathsf H\to\CP^1\) be the Hopf line and
\(\widehat c_1(\mathsf H)\in\widehat H^2(\CP^1;\ZZ)\) its differential
first Chern class.  Let
\(\widehat\omega_3\in\widehat H^3(SU(2);\ZZ)\) have curvature the
integrally normalized baryon three-form.  The correctly typed universal
object is defined on the target product, before any spacetime field is
inserted:
\begin{equation}
 \boxed{
 \widehat\Xi_5
 =2\,\operatorname{pr}_U^*\widehat\omega_3
 \smile\operatorname{pr}_s^*\widehat c_1(\mathsf H)
 \in\widehat H^5(SU(2)\times\CP^1;\ZZ).}
 \label{eq:differential-character}
\end{equation}
This is an intrinsic differential-cohomology construction, not a choice of a
five-dimensional extension~\cite{DavighiGripaiosRandalWilliams2022APE6}.
The coefficient two is the charged-two-colour mixed-WZW level matched in the
intermediate UV candidate~\cite{DavighiLohitsiri2024APE6}; global-form
refinements of symplectic QCD remain relevant~\cite{Saito2024APE6}.

Let \(\cC_{\rm restr}\) denote the smooth nonvanishing-Higgs field space
before quotienting by the gauge group.  Its evaluation map is
\begin{equation}
 \operatorname{ev}:\Sigma_3\times\cC_{\rm restr}
 \longrightarrow SU(2)\times\CP^1,\qquad
 (x,c)\longmapsto\bigl(U_c(x),s_c(x)\bigr).
 \label{eq:evaluation-map}
\end{equation}
Only now do we pull back and fiber-integrate:
\begin{equation}
 \widehat\cL_{\rm loc}
 :=\int_{\Sigma_3}\operatorname{ev}^*\widehat\Xi_5
 \in\widehat H^2(\cC_{\rm restr};\ZZ).
 \label{eq:transgression}
\end{equation}
Thus \(\widehat\cL_{\rm loc}\) is a well-typed line-with-connection
candidate on the unquotiented restricted field space.  Equation
\eqref{eq:transgression} alone does not construct an equivariant differential
cocycle or prove that the line descends to the gauge quotient stack.

\subsection{Local curvature control is not stack descent}

On a local Hopf lift define
\begin{equation}
 a=-\ii z^\dagger\dd z,
 \qquad z\mapsto e^{\ii\lambda}z,
 \qquad a\mapsto a+\dd\lambda.
 \label{eq:hopf-connection}
\end{equation}
The curvature can be written without the lift,
\begin{equation}
 f=\dd a=-\ii\,\dd z^\dagger(1-zz^\dagger)\wedge\dd z
 =s^*F_{\rm FS}.
 \label{eq:basic-curvature}
\end{equation}
It is invariant and horizontal under the vertical \(U(1)\) vector.  Since
\(U\) is neutral under vector \(U(1)_g\), the curvature
\(2\omega_3(U)\wedge f/(2\pi)\) passes the local horizontal-curvature test.
The connection \(a\) itself
is not a basic one-form; rather, its transformation law is exactly the
descent datum of the associated line bundle.  Confusing these two statements
would incorrectly demand a globally defined potential on \(\CP^1\).
Conversely, horizontal curvature does not by itself prove equivariant
differential-character descent: one must still specify the cocycle on gauge
arrows and prove trivial vertical holonomy.

\subsection{A sampled constant-loop integrality control}

For the restricted constant-loop ansatz and a gauge function with winding
\(\lambda(1)-\lambda(0)=2\pi n\), the worldline connection at baryon number
\(B\) shifts by
\begin{equation}
 \Delta S/\hbar=2B\,[\lambda(1)-\lambda(0)]=4\pi Bn.
 \label{eq:large-gauge}
\end{equation}
Therefore
\begin{equation}
 \exp(\ii\Delta S/\hbar)=1
 \qquad(B,n\in\ZZ).
 \label{eq:large-invariance}
\end{equation}
The quantized coefficient is essential.  This computation samples a
one-parameter class; it does not classify all connected components of
\(\operatorname{Map}(\Sigma_3,U(1))\), nor does it supply their equivariant
cocycle data.

\subsection{Free Hopf action and faithful centre}

The \(U(1)\) action on \(S^3\) is free whenever \(\phi\ne0\), so the
restricted Hopf quotient has no gauge stabilizer.  Separately, the faithful
global identification equates the flavour centre \(z\mapsto-z\) with the
gauge transformation \(e^{\ii\pi}\).  An \(\cO(k)\) fiber transforms by
\begin{equation}
 e^{\ii k\pi}=(-1)^k.
 \label{eq:center-holonomy}
\end{equation}
Thus \(k=2B\) has trivial phase for this tested centre element, while the
\(k=1\) negative control does not.  This is a necessary even-weight check.
It is \emph{not} an exactly-two theorem: it neither excludes higher even
weights nor checks every stabilizer, gauge component, global form, or anomaly
constraint.

\subsection{Restriction to the two candidate collective spaces}

If one formally inserts a product family
\((U_B(\bm x),s(m))\) with \(m\in\CP^1\), then
\begin{equation}
 c_1(\widehat\cL_B)
 =\int_{\Sigma_3\times\CP^1}
 2\,\omega_3(U_B)\wedge s^*\omega_2
 =2B\deg(s).
 \label{eq:formal-c1}
\end{equation}
For \(B=1,\deg(s)=1\), this is \(+2\).  But this \(s(m)\) is the
volume-divergent scalar vacuum coordinate of Eq.~\eqref{eq:bulk-norm}.
It is not a localized moduli-space map.

On the actual rotational orbit \(SO(3)\), hold the scalar vacuum fixed.  The
evaluation map has no scalar tangent component, so the two-form curvature of
Eq.~\eqref{eq:transgression} restricts trivially along the \(SO(3)\) orbit.
The WZW line therefore does not supply a free \(c_1=2\) on the physical
orientation family.

\subsection{Why even restricted stack descent remains open}

Equations~\eqref{eq:differential-character}--\eqref{eq:transgression} are
correctly typed on the unquotiented nonvanishing-Higgs field space.  To
descend them to the restricted gauge stack one must still:
\begin{enumerate}[label=(R\arabic*)]
\item treat every component of \(\operatorname{Map}(\Sigma_3,U(1))\), not
  only the sampled constant-loop ansatz;
\item construct the equivariant differential cocycle on gauge arrows and
  verify its composition law;
\item prove trivial vertical holonomy for every gauge loop.
\end{enumerate}
The extension to the full microscopic space additionally requires:
\begin{enumerate}[label=(D\arabic*)]
\item configurations where \(\phi=0\) and the projective map \(s\) fails;
\item defects and nontrivial gauge bundles not represented by one global
  normalized doublet;
\item the full faithful global-form quotient and all non-extendible bordism
  sectors of the microscopic theory;
\item a same-soliton normalizable collective map.
\end{enumerate}
Accordingly the machine card distinguishes the two flags
\begin{center}
\footnotesize
\begin{tabular}{c}
\texttt{restricted\_local\_differential\_character\_ansatz\_typed=true}\\
\texttt{restricted\_higgsed\_sigma\_model\_descent\_proven=false}\\
\texttt{gauge\_basic\_wzw\_descent\_proven=false}
\end{tabular}
\end{center}

\section{Numerical controls}

\subsection{Same-profile and stabilizer regressions}

The verifier imports the Workline-B public solver, calls it at
\(R=16\) with (4097) samples, and hard-asserts equality of the canonical
\((r,F,s)\) checksum with the Workline-B card.  It then reads \(B\) from that
same card, evaluates the Gram matrix~\eqref{eq:stabilizer-gram}, and fits the
uniform scalar-mode norm against box size.  The expected power is three.

\subsection{Finite Wilson--Dirac endpoint control}

On a cubic grid the verifier discretizes Eq.~\eqref{eq:static-hamiltonian}
with central derivatives and a Wilson control term
\begin{equation}
 H_{W,a}=H_{F,a}-\frac{r_Wa}{2}\,\beta\Delta_a.
 \label{eq:wilson-control}
\end{equation}
The mass term uses only the canonical Workline-B (F) samples, with a stated
piecewise-linear off-grid interpolation.  The accompanying canonical
breathing field (s) does not enter because the declared coupling
\eqref{eq:4d-yukawa} contains no (s)-dependent Yukawa term.  Dirichlet
finite-box boundaries are used.  The calculation checks
Hermiticity and the smallest-magnitude eigenvalues on two grids.  A positive
finite-box minimum excludes an accidental zero only for this benchmark.  It
is not a continuum extrapolation, a fermion determinant, or lattice QC2D.
The grids are deliberately small and the solver uses shift-invert
\texttt{eigsh} with a smallest-magnitude fallback for reproducibility.

There is also a direct coset-representative test.  If one incorrectly
parametrizes the hedgehog by \(SU(2)/U(1)\), then \(A\) and
\(Ae^{-\ii\theta\tau_3/2}\) denote the same proposed \(\CP^1\) point.  The
two conjugated hedgehog operators differ at nonzero \(\theta\), although they
are unitarily isospectral.  Hence the operator does not descend to that
\(\CP^1\) quotient.

\subsection{Berry, CPT, and holonomy controls}

The following controls are deliberately finite and diagnostic:
\begin{enumerate}[label=(N\arabic*)]
\item Bargmann-triangle Berry phases of Eq.~\eqref{eq:veronese} give
  \(c_1=+2\), with unit nonzero projector gap.
\item The spectral product~\eqref{eq:spectral-regulator} and its antisymmetric
  doubling obey determinant and Pfaffian conjugation to floating-point
  precision.
\item Sampled constant-loop \(U(1)\) windings give unit phase at level two;
  the faithful centre gives \(+1\) at level two and \(-1\) in the level-one
  control.  These do not prove restricted gauge-stack descent.
\end{enumerate}

\subsection{Recorded numerical values and scope}

The profile and its producing artifacts are cryptographically bound as
follows.
\begin{table}[ht]
\centering
\scriptsize
\begin{tabular}{@{}L{0.24\textwidth}L{0.69\textwidth}@{}}
\toprule
Artifact & SHA-256\\
\midrule
Canonical little-endian float64 \((r,F,s)\)
 & \nolinkurl{81104f59a5f3f4337739fc2a217cafc8da91581c9f1fb40b4fdba6ce0f948d59}\\
Workline-B solver script
 & \nolinkurl{393a1a1948b65bceb3c04167e452140a643a482a5035ed419d892fdd9cc1f10d}\\
Workline-B machine card
 & \nolinkurl{204dc3c91d3b83031fb70cafefa63e29df56f14c76c230600497ec94e1fb18b5}\\
\bottomrule
\end{tabular}
\caption{Canonical same-soliton provenance used by the C workline.}
\label{tab:provenance}
\end{table}

Table~\ref{tab:numerics} records the deterministic run.  The two Wilson
entries are deliberately reported separately: their increase with resolution
is not presented as a continuum extrapolation.  They only show that neither
finite matrix has an accidental zero at the declared benchmark.

\begin{longtable}{@{}L{0.43\textwidth}L{0.22\textwidth}L{0.27\textwidth}@{}}
\caption{AP-E6 deterministic numerical controls.}
\label{tab:numerics}\\
\toprule
Quantity & Value & Interpretation\\
\midrule
\endfirsthead
\toprule
Quantity & Value & Interpretation\\
\midrule
\endhead
Canonical Workline-B baryon number
 & \(0.999999999686391\)
 & Card-bound \(R=16\) coupled profile.\\
Canonical breathing field at the origin
 & \(s(0)=0.917128964607058\)
 & Imported and explicitly decoupled by the declared Yukawa.\\
Stabilizer Gram spectrum
 & \((16,16,16)\)
 & No continuous pure-internal stabilizer.\\
Uniform scalar norm power
 & \(3.000000000000001\)
 & Volume divergence in a cubic box.\\
Wilson \(5^3\), matrix dimension \(1000\)
 & \(\min|\lambda|=0.010896961\)
 & Finite-box endpoint control only.\\
Wilson \(7^3\), matrix dimension \(2744\)
 & \(\min|\lambda|=0.283674971\)
 & Finite-box endpoint control only.\\
False-\(\CP^1\) representative operator difference
 & \(1.15817\times10^{-2}\)
 & Relative Frobenius norm.\\
Representative isospectral residual
 & \(4.86\times10^{-17}\)
 & Unitary-conjugation check.\\
Veronese spectator Berry number
 & \(1.999999999999986\)
 & \(\cO(2)\) template, not same-model evidence.\\
Veronese kernel residual / nonzero gap
 & \(4.16\times10^{-16}/1\)
 & Projector regression.\\
CPT determinant / Pfaffian residual
 & \(0/0\)
 & Finite spectral-cutoff algebra.\\
Maximum level-two large-gauge phase residual
 & \(2.45\times10^{-15}\)
 & Sampled windings only; not stack descent.\\
\bottomrule
\end{longtable}

\begin{gate}[Meaning of a green verifier]
A green run means that the no-go deductions, conditional index formulae,
finite discretizations, and fail-closed flags are mutually consistent.  It
does not mean that a physical Callias family or portal has been constructed.
\end{gate}

\section{Gate ledger and next mathematically decisive experiments}

\begin{longtable}{@{}L{0.39\textwidth}L{0.14\textwidth}L{0.39\textwidth}@{}}
\toprule
Gate & Status & Reason\\
\midrule
\endfirsthead
\toprule
Gate & Status & Reason\\
\midrule
\endhead
Same physical \(\CP^1\) moduli space derived
 & \status{false}
 & Hedgehog orbit is \(SO(3)\); scalar \(\CP^1\) norm diverges with volume.\\
Actual Yukawa--Callias operator derived from mother model
 & \status{false}
 & Eq.~\eqref{eq:4d-yukawa} is a declared constituent completion.\\
Uniform essential Fredholm gap
 & \status{true}
 & Under (C1)--(C4), \(\Phi_\infty=\beta M\) gives essential threshold \( |M| \).\\
Uniform full endpoint spectral gap
 & \status{false}
 & Only two finite Wilson controls; no continuum exclusion of zero crossings.\\
Static ordinary Callias index \(+1\)
 & \status{false}
 & Actual asymptotic positive eigenbundle is trivial and has index zero.\\
Four-dimensional interpolation formula
 & \status{conditional theorem}
 & APS index equals spectral flow; anomaly matching gives \(B\) under stated regulator assumptions.\\
Same-model determinant line \(\cO(2)\)
 & \status{false}
 & Wrong base and no mixed mass; minimal globally gapped two-band mass is obstructed.\\
Physical CPT/anti-canonical regulator
 & \status{false}
 & Conditional algebra holds under (T1)--(T4), but no physical regulator or \(\CP^1\) polarization is derived.\\
Restricted local differential-character ansatz typed
 & \status{true}
 & Universal target character, evaluation pullback, and fiber-integration degrees are consistent.\\
Restricted Higgsed gauge-stack descent
 & \status{false}
 & All gauge-group components, an equivariant cocycle, and vertical holonomy remain unchecked.\\
Full microscopic gauge-basic descent
 & \status{false}
 & Higgs zeros, defects, all bundles/bordism sectors, and same-soliton pullback remain open.\\
Same-soliton composition / lane closed
 & \status{false}
 & No physical \(\CP^1\), Callias \(\cO(2)\), or CPT line.\\
Degree-one Route-E portal
 & \status{not authorized}
 & No pre-portal lane is closed.\\
\bottomrule
\end{longtable}

There are three concrete ways forward, ordered by how directly they address
the no-go.

\begin{enumerate}[label=\textbf{C\arabic*.}]
\item \textbf{Change the soliton, not its name.}  Search the full relaxed
  gauge--meson--scalar equations for a localized scalar texture whose exact
  symmetry breaking is \(G/H=SU(2)/U(1)\), and verify the horizontal
  \(L^2\) norm and projected Hessian gap.  This is the only route that keeps
  the desired physical \(\CP^1\) literal.
\item \textbf{Accept \(SO(3)\) and reformulate the target.}  Compute the
  torsion determinant line, Finkelstein--Rubinstein sign, and WZW holonomy on
  \(SO(3)\) or its \(SU(2)\) cover.  This may yield a meaningful two-valued
  structure, but it cannot reproduce a free Chern number two.
\item \textbf{Construct a genuine higher-rank asymptotic mass.}  Starting
  from anomaly-free microscopic fields, build a globally gapped rank-three
  or larger mass over \(S^2_\infty\times\CP^1\), compute its family index and
  heavy-regulator threshold, and show that its light kernel couples to the
  same soliton.  The Veronese control supplies only the \(q=2\) factor; the
  \(p=1\) spatial factor and decoupling are still missing.
\end{enumerate}

No degree-one portal should be started until one of these routes closes all
of its physical gates.

\bibliographystyle{unsrt}
\bibliography{route_f/tex/ap_e6_same_soliton_yukawa_callias}

\end{document}
